\subsubsection{重整化 Jensen 缺陷与排斥半径证书}\label{subsubsec:typed-address-biaxial-completion-renormalized-jensen-defect-certificate}
前一段已经说明，\(\Gamma\)-重整化把端点层的最大模账本降为多项式级；然而固定 Taylor 证书仍依赖边界最小模
\[
\widetilde m_r=\min_{|z|=1}|\widetilde F_a(rz)|,
\]
而该量对单个近边界近零点极其敏感。本段改用圆周平均的对数模长缺陷，把点态最小值替换为零点计数的积分账本。

沿用
\[
s_a(w)=\frac12+a\frac{1+w}{1-w},
\qquad
\Phi_a(w):=\widetilde F_a(w)
=
\frac12s_a(w)(s_a(w)-1)\zeta(s_a(w)).
\]
Cayley 变换给出
\[
|w|<1
\iff
\Re s_a(w)>\frac12,
\qquad
w=\frac{s-\frac12-a}{s-\frac12+a}.
\]
因子 \(s\) 在该半平面内不为零，且 \((s-1)\zeta(s)\) 在 \(s=1\) 处可去并取值 \(1\)。因此 \(\Phi_a\) 在单位圆盘内的零点恰为 \(\zeta(s)\) 在 \(\Re s>1/2\) 内零点的 Cayley 像。于是
\[
\mathrm{RH}
\iff
\Phi_a(w)\ne0\qquad(|w|<1).
\]

\begin{lemma}[重整化基点非零]\label{lem:typed-address-biaxial-completion-renormalized-basepoint-nonzero}
对任意 \(a>0\)，有
\[
\Phi_a(0)\ne0.
\]
\end{lemma}
\begin{proof}
记 \(\sigma_0=s_a(0)=1/2+a\)。若 \(\sigma_0>1\)，则 Euler 乘积给出 \(\zeta(\sigma_0)\ne0\)。若 \(\sigma_0=1\)，则
\[
\frac12s(s-1)\zeta(s)\longrightarrow \frac12
\qquad(s\to1).
\]
若 \(1/2<\sigma_0<1\)，则
\[
\zeta(\sigma_0)=\frac{\eta(\sigma_0)}{1-2^{1-\sigma_0}},
\qquad
\eta(\sigma)=\sum_{n\ge1}\frac{(-1)^{n-1}}{n^\sigma}.
\]
交错级数判别给出 \(\eta(\sigma_0)>0\)，而 \(1-2^{1-\sigma_0}\ne0\)，故 \(\zeta(\sigma_0)\ne0\)。三种情形均推出 \(\Phi_a(0)\ne0\)。证毕。
\end{proof}

\begin{definition}[重整化 Jensen 缺陷]\label{def:typed-address-biaxial-completion-renormalized-jensen-defect}
对 \(0<r<1\)，定义
\[
J_a(r)
:=
\frac1{2\pi}
\int_0^{2\pi}\log|\Phi_a(re^{i\theta})|\,\dd\theta
-\log|\Phi_a(0)|.
\]
\end{definition}

\begin{proposition}[Jensen 缺陷的零点计数账本]\label{prop:typed-address-biaxial-completion-renormalized-jensen-zero-count-ledger}
设 \((\alpha)\) 遍历 \(\Phi_a\) 在 \(\DD\) 内的零点，按重数 \(m_\alpha\) 计数。则
\[
J_a(r)
=
\sum_{|\alpha|<r}m_\alpha\log\frac r{|\alpha|}.
\]
特别地，
\[
J_a(r)\ge0,
\qquad
J_a(r_2)-J_a(r_1)
=
\int_{r_1}^{r_2}N_a(t)\frac{\dd t}{t}
\quad(0<r_1<r_2<1),
\]
其中
\[
N_a(t):=\sum_{|\alpha|<t}m_\alpha.
\]
在没有零点落在 \(|w|=r\) 的半径处，
\[
rJ_a'(r)=N_a(r).
\]
\end{proposition}
\begin{proof}
由引理 \ref{lem:typed-address-biaxial-completion-renormalized-basepoint-nonzero}，Jensen 公式可用于 \(\Phi_a\)，并给出所示求和表达；若半径经过零点模长，则以径向极限理解，公式仍保持同一阶梯计数语义。将每个零点项写成
\[
\log\frac r{|\alpha|}
=
\int_{|\alpha|}^{r}\frac{\dd t}{t}
\qquad(|\alpha|<r),
\]
并交换有限区间内的离散求和与积分，即得增量公式。对 \(\log r\) 求导则得到 \(rJ_a'(r)=N_a(r)\)。证毕。
\end{proof}

\begin{theorem}[RH 的重整化 Jensen 等价表述]\label{thm:typed-address-biaxial-completion-renormalized-jensen-rh-equivalence}
下列命题等价：
\[
\mathrm{RH},
\qquad
J_a(r)=0\quad(0<r<1),
\qquad
\lim_{r\uparrow1}J_a(r)=0.
\]
等价地，对任意共终半径序列 \(r_q\uparrow1\)，
\[
\mathrm{RH}
\iff
J_a(r_q)\longrightarrow0.
\]
\end{theorem}
\begin{proof}
RH 等价于 \(\Phi_a\) 在 \(\DD\) 内无零点；此时命题 \ref{prop:typed-address-biaxial-completion-renormalized-jensen-zero-count-ledger} 的求和为空，故 \(J_a(r)=0\) 对所有 \(r<1\) 成立。反之，若存在零点 \(\alpha\in\DD\)，则任取 \(r>|\alpha|\) 有
\[
J_a(r)\ge m_\alpha\log\frac r{|\alpha|}>0,
\]
并且
\[
\liminf_{r\uparrow1}J_a(r)
\ge
m_\alpha\log\frac1{|\alpha|}>0.
\]
故极限趋零与恒零均排除盘内零点。共终序列形式由 \(J_a\) 的单调性立即得到。证毕。
\end{proof}

\begin{proposition}[Jensen 排斥半径证书]\label{prop:typed-address-biaxial-completion-renormalized-jensen-repulsion-certificate}
设 \(0<R<r<1\)。则闭盘计数函数
\[
N_a^{\le}(R):=\sum_{|\alpha|\le R}m_\alpha
\]
满足
\[
N_a^{\le}(R)\log\frac rR\le J_a(r).
\]
因此
\[
J_a(r)<\log\frac rR
\quad\Longrightarrow\quad
\Phi_a(w)\ne0\qquad(|w|\le R).
\]
更一般地，若 \(J_a(r)\le\delta\)，则
\[
\Phi_a(w)\ne0\qquad(|w|<re^{-\delta}).
\]
\end{proposition}
\begin{proof}
每个满足 \(|\alpha|\le R\) 的零点对 \(J_a(r)\) 的贡献至少为 \(\log(r/R)\)，故第一式成立。若 \(J_a(r)<\log(r/R)\)，则整数 \(N_a^{\le}(R)\) 必为 \(0\)。若 \(J_a(r)\le\delta\)，而某零点满足 \(|\alpha|<re^{-\delta}\)，则
\[
J_a(r)\ge\log\frac r{|\alpha|}>\delta,
\]
矛盾。证毕。
\end{proof}

\begin{corollary}[dyadic 半径层的 Jensen 排除证书]\label{cor:typed-address-biaxial-completion-renormalized-jensen-dyadic-certificate}
固定 \(0<\eta<1\)，并令
\[
R_q:=1-2^{-q},
\qquad
r_q^+:=1-\eta2^{-q}.
\]
若
\[
J_a(r_q^+)
<
\log\frac{r_q^+}{R_q},
\]
则
\[
\Phi_a(w)\ne0\qquad(|w|\le R_q).
\]
进一步地，对任意固定 \(0<\eta<1\)，RH 等价于上述不等式对所有充分大的 \(q\) 成立。
\end{corollary}
\begin{proof}
第一句是命题 \ref{prop:typed-address-biaxial-completion-renormalized-jensen-repulsion-certificate} 在 \(R=R_q\)、\(r=r_q^+\) 下的直接代入。若 RH 成立，则 \(J_a(r)=0\)，故不等式成立。反过来，若该不等式对所有充分大的 \(q\) 成立，则每个闭盘 \(|w|\le1-2^{-q}\) 在充分大 \(q\) 上无零点；这些闭盘共终穷尽 \(\DD\)，故 \(\Phi_a\) 在 \(\DD\) 内无零点，从而 RH 成立。证毕。
\end{proof}

该 dyadic 阈值具有端点尺度
\[
\log\frac{1-\eta2^{-q}}{1-2^{-q}}
=(1-\eta)2^{-q}+O(2^{-2q}).
\]
因此，排除 \(|w|\le1-2^{-q}\) 内零点所需的 Jensen 缺陷上界必须达到 \(2^{-q}\) 同阶。与边界最小模不同，\(J_a(r)\) 是非负、单调且具有对数半径导数计数律的平均型对象；它把近边界近零点的影响压成可审计的缺陷质量，而不是让单点最小值支配整个证书。

\begin{proposition}[单个离线零点的 Jensen 印记]\label{prop:typed-address-biaxial-completion-renormalized-jensen-single-offline-imprint}
若 \(\rho=\beta+i\gamma\)、\(\beta>1/2\)，并记 \(\delta=\beta-1/2>0\)，则其圆盘像
\[
\alpha_\rho
=
\frac{\delta-a+i\gamma}{\delta+a+i\gamma}
\]
满足
\[
1-|\alpha_\rho|^2
=
\frac{4a\delta}{(\delta+a)^2+\gamma^2},
\qquad
\log\frac1{|\alpha_\rho|}
=
\frac12\log
\frac{(\delta+a)^2+\gamma^2}{(\delta-a)^2+\gamma^2}.
\]
当 \(|\gamma|\to\infty\) 时，
\[
1-|\alpha_\rho|\sim \frac{2a\delta}{\gamma^2},
\qquad
\log\frac1{|\alpha_\rho|}
\sim
\frac{2a\delta}{\gamma^2}.
\]
\end{proposition}
\begin{proof}
由 \(s_a(w)=\rho\) 反解得到所示 \(\alpha_\rho\)。取模平方即得 \(1-|\alpha_\rho|^2\) 的闭式；再由
\[
\log\frac1{|\alpha_\rho|}
=
\frac12\log\frac{(\delta+a)^2+\gamma^2}{(\delta-a)^2+\gamma^2}
\]
并展开 \(|\gamma|\to\infty\) 的主项，得到渐近式。证毕。
\end{proof}

命题 \ref{prop:typed-address-biaxial-completion-renormalized-jensen-single-offline-imprint} 表明，高度 \(|\gamma|\) 的离线零点在 Jensen 缺陷边界极限中留下 \(O(\gamma^{-2})\) 级质量。这与固定 Cayley 图表中的端点二次压缩完全一致。由此，证书瓶颈从点态最小模转化为平均缺陷的同尺度上界问题：若要直接排除 dyadic 层 \(|w|\le1-2^{-q}\)，必须将
\[
J_a(1-\eta2^{-q})
\]
严格压到 \((1-\eta)2^{-q}+O(2^{-2q})\) 以下。后续的有限积分、Fourier 或 Toeplitz--PSD 编译，正是对这一平均缺陷上界的离线证书化。
